\chapter{Search Methods}
\label{chap:algorithms}

% Sources: src/hrgs_scheduler/search/, docs/Optimizer Status.md,
% docs/Justification of Implementation.md §4

Three search tiers share one schedule representation and one evaluator. Every returned candidate passes the structural checks of Section~\ref{sec:model:legality} and is evaluated by the bottom-up pass of Section~\ref{sec:model:objectives}; the tiers differ only in which legal \Glspl{dag} they consider, never in how a returned \Gls{dag}'s fidelity, success probability, latency, or cost is calculated.

Section~\ref{sec:alg:prelims} introduces the general concepts used throughout the search, followed by the shared evaluator in Section~\ref{sec:alg:evaluator}. The search is organised into four tiers with increasing scope but decreasing coverage of the schedule space: fixed baseline families in Section~\ref{sec:alg:bruteforce}, the exact span-partition recurrence in Section~\ref{sec:alg:dp}, same-span pumping in Section~\ref{sec:alg:pumping}, and beam search in Section~\ref{sec:alg:beam}. The latter is the heuristic used for the production-scale experiments in Chapter~\ref{chap:experiments}.

\section{Preliminaries: Combinatorial Search Concepts}
\label{sec:alg:prelims}

The schedule space defined in Section~\ref{sec:model:dag} is combinatorially large: even modest chains admit too many legal schedules to enumerate individually. The three tiers below rely on three standard tools from combinatorial optimisation, introduced briefly here before being applied to schedules.

\vspace{0.5cm}

\paragraph{Multi-objective optimisation and Pareto dominance} A search problem is \emph{multi-objective} when candidates are compared on several quantities at once, here cost $C$, fidelity $F$, and success probability $P_{\mathrm{success}}$ (Section \ref{sec:model:objectives}). Candidate $A$ \emph{Pareto-dominates} $B$ if $A$ is at least as good on every quantity and strictly better on one; a \emph{Pareto-optimal} candidate is dominated by none, and the \emph{Pareto frontier} is the set of all such undominated candidates~\cite{Pareto2014Manual,Deb2002Fast,Ehrgott2005Multicriteria}. Retaining the whole frontier matters here because a locally inferior candidate can become the better choice once combined with another resource, as happens when two spans are swapped together.

\vspace{0.5cm}

\paragraph{Dynamic programming} \emph{Dynamic programming} decomposes a large problem into smaller overlapping subproblems that are solved once and reused wherever they recur (\emph{memoized})~\cite{Bellman2003Dynamic}. Section~\ref{sec:alg:dp} applies this principle to schedules: candidate sub-schedules for a short span are computed once and reused when constructing longer spans.

\vspace{0.5cm}

\paragraph{Heuristic and beam search} When candidates become too numerous to enumerate exhaustively, a \emph{heuristic search} trades completeness for tractability by discarding some before they are fully explored. \emph{Beam search} does this by keeping only a fixed number of candidates, the \emph{beam width}, at each stage, ranked by a chosen rule~\cite{Lowerre1990HARPY,Huang2018When,Meister2022BestFirst}; a narrower beam is cheaper but more likely to discard a candidate that would have led to a better final schedule. Section~\ref{sec:alg:beam} applies this to the span-partition recursion of Section~\ref{sec:alg:dp}.

\clearpage

\section{The Schedule Evaluator}
\label{sec:alg:evaluator}

% Source: schedule/evaluator.py, experiments/fig5_fidelity_vs_noise.py

The evaluator makes one bottom-up pass over a schedule \Gls{dag}. At each node it applies the corresponding physical operation to the current Pauli-channel state, synchronizes branches that become ready at different times, applies the resulting memory decoherence, and propagates side effects and timing. At the root it returns fidelity $F$, success probability $P_{\mathrm{success}}$, resource cost $C$, and latency $L$. Rate is then $R=P_{\mathrm{success}}/L$ under the full-restart model of Section~\ref{sec:model:objectives}.

Section~\ref{sec:exp:validation} shows that this pass reproduces the source paper's Fig.\,5 fidelity numbers to four significant figures, confirming that the evaluator reproduces the source model correctly.

\section{Fixed Schedule Families}
\label{sec:alg:bruteforce}

The first tier enumerates four fixed structural families:
\begin{itemize}
    \item a \emph{single raw end-to-end chain}, with no purification;
    \item \emph{end-node pumping}, in which $n_{\mathrm{pur}}$ independent raw chains are purified sequentially at the end nodes with a round-trip herald between rounds (\emph{heralded} end-node pumping);
    \item its \emph{optimistic counterpart}, deferring heralding until the end (\emph{optimistic} end-node pumping);
    \item \emph{uniform link-level} purification, using the same copy count and circuit at every hop.
\end{itemize}
For even $N$ within budget, the Integrating paper's hand-chosen schedule~\cite{Benchasattabuse2025Integrating} is included as a fifth fixed comparison point.

Three of these families, the raw chain, heralded end-node pumping, and optimistic end-node pumping, are referred to throughout this thesis as the \emph{canonical schedules}: they are the schedules whose fidelity is checked directly against the source paper's published Fig.\,5 benchmarks (Section~\ref{sec:exp:validation}).

Within each family, purification sequences use the three two-copy circuits $\mathcal C_{\mathrm{YY}}$, $\mathcal C_{\mathrm{ZX}}$, and $\mathcal C_{\mathrm{XZ}}$ (Section~\ref{sec:bg:circuits}). All $3^r$ circuit sequences are enumerated up to a configured round cutoff $r_{\max}$; above that cutoff, the implementation uses a small curated set. Thus, ``brute force'' denotes exhaustive enumeration within these predefined structural families and circuit grid, not enumeration of every legal schedule \Gls{dag}.

These fixed families serve two purposes. First, they provide explicit baselines throughout Chapter~\ref{chap:experiments}: the raw chain and source-paper schedule serve as reference points, while uniform link-level purification provides a simple non-adaptive baseline without structural search. Second, their candidates are always retained alongside those produced by Second, their candidates are always kept alongside candidates from the broader searches, so they cannot be removed by pruning.

\section{Span Dynamic Programming}
\label{sec:alg:dp}

For every span $\mathrm{Span}(a,b)$, the second tier memoizes a Pareto frontier~\cite{Pareto2014Manual} of ways to create one resource over that span. A single-hop frontier ($b=a+1$) contains the raw link and link-level purification choices. For a wider span, the recurrence tries every split point $a<m<b$ and combines each retained candidate from $\mathrm{Span}(a,m)$ with each from $\mathrm{Span}(m,b)$ using a Swap operation (Section~\ref{sec:bg:hrgs}). Each sub-span is computed once and reused in wider constructions. This avoids repeatedly constructing and evaluating complete end-to-end schedules, which is the main computational benefit of the dynamic program.

\clearpage

The frontier retained at each span is multi-objective over $(C,F,P_{\mathrm{success}})$ (Pareto dominance, Section~\ref{sec:alg:prelims}). Only non-dominated candidates are retained, following the standard Pareto-pruning rule~\cite{Ehrgott2005Multicriteria}. This pruning is exact with respect to the three quantities tracked by the frontier but does not make the search exhaustive. The caps and heuristic retention rules introduced in the next two sections can still discard a non-dominated candidate before it reaches the end-to-end span.

\vspace{0.7cm}

\paragraph{Same-Span Entanglement Pumping}
\label{sec:alg:pumping}

The span-partition recurrence described above only builds a resource by swapping two \emph{different} spans, so it cannot represent purifying a resource against a fresh, independently generated candidate at the \emph{same} span $\kappa$. The second-tier extension allows this operation with limits on the number of copies and circuits used. The uncapped mode is used only for targeted validation, as it becomes impractical even at small $N$ for large budgets (see Table~\ref{tab:benchmarks} in Appendix~\ref{app:searchalgos}).

Pumping candidates and ordinary Swap candidates share the same fixed-size frontier. As a result, adding pumping candidates can remove Swap candidates from the frontier before those candidates can be extended to a complete end-to-end schedule. Disabling pumping entirely reproduces pumping-free search behaviour; Section~\ref{sec:exp:headline} reports both settings at the reference configuration since they can return different top-ranked schedules.

\vspace{0.6cm}

\section{Beam Search}
\label{sec:alg:beam}

The third tier, beam search, reuses the same span recursion, node construction, and evaluator as Section~\ref{sec:alg:dp}, but caps every retained span frontier at a beam width $w_{\text{beam}}$, bounding runtime and memory at the cost of possibly pruning a candidate a wider beam would keep. The cap is necessary in practice: an uncapped width-7 frontier already holds roughly $1{,}200$ candidates, making an uncapped search through the paper's $N=10$ chain impractical. The same $w_{\text{beam}}$ also caps the same-span pumping pool, so pumping and ordinary swap candidates share one set of slots at each span.

At each span, roughly half of the $w_{\text{beam}}$ slots are reserved for the highest-fidelity candidates, while the rest favour candidates that meet $F_{\min}$ and have high success probability and low cost. This prevents the search from discarding link-level purification simply because it is less efficient locally. Such candidates may still be needed to maintain sufficient fidelity after several swaps.

Throughout this thesis, a \emph{score} is the returned candidate's value of whichever quantity the active objective maximizes (Section~\ref{sec:model:objectives}). Under the default rate-maximizing objective used for the beam-search benchmarks of Table~\ref{tab:benchmarks} in Appendix~\ref{app:searchalgos}, score is simply the candidate's rate $R(\Sigma)$.

A different heuristic (e.g.\ simulated annealing or a genetic search) would need its own way of representing schedules and generating new candidates while respecting the legality conditions of Section~\ref{sec:model:legality}. Beam search instead keeps the same span-partition representation and only changes how many candidates survive at each frontier, keeping every tier on one physical evaluation path and every run \emph{deterministic}. As stated before, any given result thus only represents a valid candidate found by a bounded search, not a guarantee of global optimality.